\begin{pastexam}[label={px:ccu-113-swa-09}]{中正資工所甲組 113 學年度 第 9 題}
（單選）

Given the directed graph below, please answer the following questions from Q8 to Q9.
\begin{center}
\begin{tikzpicture}[font=\sffamily\small,>={Stealth[length=2mm]}]
\node[draw,circle,minimum size=8mm,inner sep=0pt,font=\sffamily\bfseries] (A) at (1.24,-3.49) {A};
\node[draw,circle,minimum size=8mm,inner sep=0pt,font=\sffamily\bfseries] (B) at (2.53,-1.85) {B};
\node[draw,circle,minimum size=8mm,inner sep=0pt,font=\sffamily\bfseries] (F) at (5.35,-1.12) {F};
\node[draw,circle,minimum size=8mm,inner sep=0pt,font=\sffamily\bfseries] (D) at (4.29,-3.49) {D};
\node[draw,circle,minimum size=8mm,inner sep=0pt,font=\sffamily\bfseries] (E) at (7.82,-3.49) {E};
\node[draw,circle,minimum size=8mm,inner sep=0pt,font=\sffamily\bfseries] (C) at (5.06,-5.47) {C};
\draw[->] (A)--node[above left,inner sep=1pt]{10}(B);
\draw[->] (B)--node[above,inner sep=2pt]{15}(F);
\draw[->] (B)--node[below left,inner sep=1pt]{12}(D);
\draw[->] (D)--node[right,inner sep=4pt]{1}(F);
\draw[->] (D)--node[above,inner sep=2pt]{2}(E);
\draw[->] (F)--node[above right,inner sep=1pt]{5}(E);
\draw[->] (A)--node[below left,inner sep=1pt]{15}(C);
\draw[->] (C)--node[below right,inner sep=1pt]{10}(E);
\end{tikzpicture}
\end{center}

(2\%) What is the topological ordering of the graph? You must choose the node with the label that has the smallest alphabetical order first, when finding the topological order.
\begin{enumerate}[label=\Alph*.]
\item A$\to$B$\to$D$\to$F$\to$E$\to$C
\item A$\to$C$\to$E$\to$B$\to$D$\to$F
\item A$\to$B$\to$D$\to$E$\to$C$\to$F
\item A$\to$B$\to$C$\to$D$\to$F$\to$E
\item None of the above
\end{enumerate}
\end{pastexam}

\begin{pastsoln}
答案是 D（$A \to B \to C \to D \to F \to E$）。 邊上的數字是第 8 題 Dijkstra 用的成本，拓樸排序用不到。

「每次挑字母最小的」就是資結書第~22.1.3~節用最小優先佇列的 Kahn 演算法：每次從入度為 $0$ 的頂點中取字母最小的輸出，
再把它出去的邊拿掉。 一開始的入度是 $A{:}\,0$、$B{:}\,1$、$C{:}\,1$、$D{:}\,1$、$E{:}\,3$、$F{:}\,2$。 表中「--」表示已經輸出。
\begin{center}
\small
\begin{tabular}{cc*{6}{c}l}
\toprule
步驟 & 取出 & $A$ & $B$ & $C$ & $D$ & $E$ & $F$ & 可取的頂點 \\
\midrule
開始 & & 0 & 1 & 1 & 1 & 3 & 2 & $\{A\}$ \\
1 & $A$ & -- & 0 & 0 & 1 & 3 & 2 & $\{B, C\}$ \\
2 & $B$ & -- & -- & 0 & 0 & 3 & 1 & $\{C, D\}$ \\
3 & $C$ & -- & -- & -- & 0 & 2 & 1 & $\{D\}$ \\
4 & $D$ & -- & -- & -- & -- & 1 & 0 & $\{F\}$ \\
5 & $F$ & -- & -- & -- & -- & 0 & -- & $\{E\}$ \\
6 & $E$ & -- & -- & -- & -- & -- & -- & （空） \\
\bottomrule
\end{tabular}
\end{center}

其他選項：
\begin{itemize}
\item A：$A, B, D, F, E, C$ 違反邊 $C \to E$（$E$ 排在 $C$ 前面），根本不是拓樸順序。 第 3 步時 $C$、$D$ 都可以取，照規則也
該先取 $C$。
\item B：$E$ 排在 $B$、$D$、$F$ 前面，違反 $D \to E$、$F \to E$。
\item C：$E$ 排在 $C$、$F$ 前面，違反 $C \to E$、$F \to E$。
\item E：D 是對的，所以不選。
\end{itemize}
第 8、9 題的答案剛好都是 D：這張圖的 Dijkstra 取出順序本身也是一個拓樸順序（每條邊都從先取出的指向後取出的），但這是巧合，
兩個問題一般沒有關係。
\end{pastsoln}
